\documentclass[conference]{IEEEtran}
\IEEEoverridecommandlockouts
\usepackage{cite}
\usepackage{amsmath,amssymb,amsfonts}
\usepackage{algorithmic}
\usepackage{graphicx}
\usepackage{textcomp}
\usepackage{xcolor}

\def\BibTeX{{\rm B\kern-.05em{\sc i\kern-.025em b}\kern-.08em
    T\kern-.1667em\lower.7ex\hbox{E}\kern-.125emX}}

\begin{document}

\title{[PAPER TITLE — specific, describes the mechanism, not a vague theme]}

\author{
\IEEEauthorblockN{[Author Name]}
\IEEEauthorblockA{\textit{[Department]} \\
\textit{[Institution]}\\
[City, Country] \\
[email]}
}

\maketitle

\begin{abstract}
[150-250 words. Must be self-contained: problem, approach, the single strongest result with an actual number, and the takeaway. No citations inside the abstract. No unsupported superlatives — every claim of improvement here must be backed by a metric stated later in the paper.]
\end{abstract}

\begin{IEEEkeywords}
[keyword1, keyword2, keyword3, keyword4, keyword5]
\end{IEEEkeywords}

\section{Introduction}
% Problem statement, why it matters, what existing approaches miss (specific mechanism of failure, not "prior work is limited"),
% the paper's contribution as explicit bullet points, and a one-sentence roadmap of the remaining sections.

\section{Related Work}
% Position against the nearest work explicitly -- table or paragraph comparing this work's mechanism against each
% close competitor's mechanism, not just a list of citations. Run the researcher / researcher-evaluator skills on this
% section specifically before treating it as complete: an incomplete related-work section is the single most common
% reviewer rejection reason after weak results.

\section{Proposed Method}
% Enough detail for a reader to reproduce this without access to the code. Define notation before using it.
% State design decisions and, where non-obvious, why the alternative was rejected.

\section{Experimental Setup}
% Dataset(s), baselines (must include the strongest real competitor, not just an easy strawman), metrics,
% hyperparameters, hardware -- enough to reproduce.

\section{Results and Discussion}
% Report negative/limiting results too. Every improvement claim must cite the specific metric and the specific
% baseline it beat. A results section with no failure case or limitation reads as unreviewed, not as strong.

\section{Limitations}
% State explicitly what this does not solve and under what conditions it would break. Reviewers penalize
% omitted limitations far more than they penalize disclosed ones.

\section{Conclusion and Future Work}
% Restate the contribution against the original problem statement -- not a generic "we presented a novel approach."

\begin{thebibliography}{00}
\bibitem{b1} [Author initials], "[Title in quotes]," \textit{[Venue]}, [Year], pp. [pages].
\end{thebibliography}

\end{document}
